芯片热流密度不断攀升，对散热能力提出了更高要求。微通道流动沸腾利用相变潜热散热，被认为是解决高热流密度散热问题最有前景的技术之一[1-3]。然而，流动沸腾不稳定性会诱发局部干涸，导致临界热流密度显著降低，严重制约了该技术的工程应用。

申请人前期研究表明，在微通道底面构筑亲疏水相间的纳米结构可以抑制流动不稳定性：在质量流速为200 kg/(m2·s)的条件下，压降波动幅值从35 kPa降低到9 kPa，临界热流密度提高了约42%。在此基础上，本项目拟提出“润湿性梯度调控汽化核心分布”的研究思路，研究成果有望为高性能芯片散热器的设计提供理论依据。
